\section{Results and Discussion}

\subsection{Catalytic-Residue Swap Profiles in Human and Mouse GPX6}

To isolate the energetic effect of residue~49 before the pathway analysis, we first compared wild-type human Sec-GPX6 and mouse Cys-GPX6 with their reciprocal catalytic-residue variants, human Cys-GPX6 and mouse Sec-GPX6. Figure~\ref{fig:sec_vs_cys_profiles} shows the bin-averaged EVB free-energy profiles for this four-system catalytic-residue swap. Human Sec-GPX6 has an activation barrier of $16.40 \pm 0.68$~kcal/mol, whereas human Cys-GPX6 increases to $17.83 \pm 0.97$~kcal/mol. Mouse Cys-GPX6 has a barrier of $17.68 \pm 0.91$~kcal/mol, whereas mouse Sec-GPX6 shows the lowest barrier, $14.84 \pm 0.52$~kcal/mol. Thus, the figure defines the two wild types as human Sec-GPX6 and mouse Cys-GPX6 and shows that the Cys-containing profiles are comparable, while the Sec-containing profiles shift downward to an extent that depends on the homologous scaffold.

 These data show that, as expected, replacement of Cys by Sec lowers the barrier for the proton-transfer step in both sequence backgrounds. The barrier reduction is $1.43$~kcal/mol in the human scaffold and $2.84$~kcal/mol in the mouse scaffold, and in both cases the effect exceeds the combined standard errors of the profiles. This residue-49 swap therefore confirms the catalytic advantage of Sec in GPX6, while also showing that the magnitude of this advantage is modulated by the surrounding protein environment. 

\begin{figure}[H]
\centering
\begin{subfigure}{0.49\textwidth}
    \centering
    \includegraphics[width=\linewidth]{Fig3A.png}
    \caption{Human Cys-GPX6.}
    \label{fig:sec_vs_cys_human_cys}
\end{subfigure}
\hfill
\begin{subfigure}{0.49\textwidth}
    \centering
    \includegraphics[width=\linewidth]{Fig3B.png}
    \caption{Human Sec-GPX6 (wt).}
    \label{fig:sec_vs_cys_human_sec}
\end{subfigure}
\vspace{0.5em}
\begin{subfigure}{0.49\textwidth}
    \centering
    \includegraphics[width=\linewidth]{Fig3C.png}
    \caption{Mouse Cys-GPX6 (wt).}
    \label{fig:sec_vs_cys_mouse_cys}
\end{subfigure}
\hfill
\begin{subfigure}{0.49\textwidth}
    \centering
    \includegraphics[width=\linewidth]{Fig3D.png}
    \caption{Mouse Sec-GPX6.}
    \label{fig:sec_vs_cys_mouse_sec}
\end{subfigure}
\caption[Catalytic-residue swap EVB profiles]{EVB free-energy profiles for the residue-49 catalytic swap in human and mouse GPX6. Each panel shows bin-averaged profiles from independent replicas for the proton-transfer step. Human Sec-GPX6 and mouse Cys-GPX6 are the wild-type systems; human Cys-GPX6 and mouse Sec-GPX6 are the reciprocal catalytic-residue variants.}
\label{fig:sec_vs_cys_profiles}
\end{figure}

\subsection{Mutational Landscapes and Pathway Asymmetry}

We then examined how the 20 nearby substitutions alter the two proteins after the residue-49 exchange. The mouse-to-human path starts from mouse Sec-GPX6 (C49U) and adds human residues, whereas the human-to-mouse path starts from human Cys-GPX6 (U49C) and adds mouse residues. At each step, we retained the substitution with the lowest calculated activation barrier. Figure~\ref{SI-fig:gpx6_all_pathways} shows this selected path together with all alternatives that were calculated. The first noncatalytic substitution is at residue~54 in the mouse-to-human path and at residue~87 in the human-to-mouse path. Figure~\ref{fig:pathway_heatmaps} summarizes the cumulative energetic effects. Each column is one step of the path, and each row shows how a substitution contributes to $\Delta G^{\ddagger}$ and $\Delta G^{\circ}$ relative to the C49U or U49C starting variant.

The heatmaps show a clear difference between directions. When human residues are added to mouse Sec-GPX6, most steps lower the activation barrier or leave it nearly unchanged. By contrast, several of the first mouse substitutions added to human Cys-GPX6 raise the barrier. Thus, the effect of a substitution depends on whether it is introduced into the C49U mouse protein or the U49C human protein. Reaction free energies also differ between directions, although the contrast is strongest for the activation barriers. Both paths remain accessible, but the same sequence changes do not have equivalent effects in the two proteins.

This dependence on sequence context is consistent with epistasis. Substitutions that are neutral or favorable with either Sec or Cys at position~49 could, in principle, have preceded the loss of Sec. Substitutions whose effect changes markedly with residue~49 are more likely to have favored one catalytic state over the other. The EVB calculations do not reconstruct the historical order of mutations, but they identify substitutions whose energetic effect depends on the presence of Sec or Cys.

\begin{figure}[H]
    \centering
    \begin{subfigure}{\textwidth}
        \centering
        \includegraphics[width=0.95\textwidth]{Fig4A.png}
        \caption{Mouse-to-human pathway.}
        \label{fig:heatmap_mouse}
    \end{subfigure}
    \vspace{0.5em}
    \begin{subfigure}{\textwidth}
        \centering
        \includegraphics[width=0.95\textwidth]{Fig4B.png}
        \caption{Human-to-mouse pathway.}
        \label{fig:heatmap_human}
    \end{subfigure}
    \caption[Cumulative free-energy changes along greedy pathways]{Cumulative free-energy changes along the selected GPX6 substitution paths. Heatmaps show changes in activation free energy ($\Delta G^{\ddagger}$) and reaction free energy ($\Delta G^{\circ}$) as substitutions are added in the mouse-to-human and human-to-mouse directions. Columns correspond to successive steps, rows correspond to individual substitutions, negative values indicate stabilization, and positive values indicate destabilization relative to the C49U mouse or U49C human starting variant; colors denote distance from catalytic residue~49.}
    \label{fig:pathway_heatmaps}
\end{figure}

\begin{figure}[H]
    \centering
    \begin{subfigure}{\textwidth}
        \centering
        \includegraphics[width=0.88\textwidth,height=0.34\textheight,keepaspectratio]{Fig5A.png}
        \caption{Mouse GPX6 variants.}
        \label{fig:lfer_mouse}
    \end{subfigure}
    \vspace{0.5em}
    \begin{subfigure}{\textwidth}
        \centering
        \includegraphics[width=0.88\textwidth,height=0.34\textheight,keepaspectratio]{Fig5B.png}
        \caption{Human GPX6 variants.}
        \label{fig:lfer_human}
    \end{subfigure}
    \caption[Linear free-energy relationships]{Linear free-energy relationships for mouse and human GPX6 variants. Each panel correlates the change in reaction free energy ($\Delta G^{\circ}$) with the change in activation free energy ($\Delta G^{\ddagger}$). Grey points represent individual simulation replicas, colored points show mutation averages with standard errors, colors denote distance from catalytic residue~49, and the red line shows the linear regression fit. The correlations are moderate in both homologs ($R^2 = 0.51$ for mouse and $R^2 = 0.54$ for human), indicating that thermodynamic perturbations are partly, but not completely, propagated into the catalytic barrier. }
    \label{fig:lfer_combined}
\end{figure}

The linear free-energy relationship analysis in Figure~\ref{fig:lfer_combined} shows that thermodynamic--kinetic coupling differs only modestly between the two homologs. Substitutions that destabilize the reaction also tend to increase the activation barrier in both datasets, with $R^2 = 0.51$ for mouse GPX6 and $R^2 = 0.54$ for human GPX6. These moderate correlations indicate that changes in reaction thermodynamics are partly propagated into changes in the activation barrier. At the residue level, however, the energetic effects can change sign across mutational levels, pointing to contributions beyond local hydrogen-bonding in the full energetic response (Figure~\ref{SI-fig:heatmap_levels}).

Figure~\ref{fig:gpx6_human_to_mouse_profile} shows the activation barriers along the human-to-mouse path, beginning with U49C. The barrier does not decrease monotonically, but the changes are smaller than the spread among all calculated alternatives (Figure~\ref{SI-fig:gpx6_all_pathways}). In both directions, the selected path remains close to the lowest barriers found at each step, whereas many alternatives give substantially higher values.

\begin{figure}[H]
    \centering
    \includegraphics[width=\textwidth]{Fig6.png}
    \caption[Human-to-mouse activation free-energy profile]{Activation free-energy ($\Delta G^{\ddagger}$) profile along the human-to-mouse mutational pathway when the U49C mutation is introduced at each mutation level. Error bars indicate standard deviation.}
    \label{fig:gpx6_human_to_mouse_profile}
\end{figure}

\subsection{Phylogenetic Context of the EVB Pathways}

Figure~\ref{fig:gpx6_triangle} places the EVB-derived trajectories in the phylogenetic framework of Rees et al.\cite{rees_ancient_2024} The ancestral reference is the reconstructed Sec-containing ancestor before the primate and rodent/lagomorph GPX6 lineages diverged, referred to as node~25 in that study.

Variants from the mouse-to-human trajectory and variants from the human-to-mouse trajectory form largely distinct clusters in the 20-residue sequence-distance analysis, while the reconstructed ancestor behaves as a transitional reference rather than as a recovered endpoint of either greedy path (Figure~\ref{SI-fig:gpx6_heatmap}). This supports the interpretation that the EVB paths capture energetically accessible transformations between modern homologs and place those transformations in relation to the reconstructed ancestral sequence.

The geometric projection in Figure~\ref{fig:gpx6_triangle} maps individual intermediates relative to mouse wild-type GPX6, human wild-type GPX6, and the node~25 ancestral sequence. The trajectories do not collapse onto a single line, and the endpoints do not perfectly exchange positions in this reduced projection. Instead, both directional pathways move through partially overlapping regions near the base of the triangle while retaining lineage-specific structure. This behavior argues against a strict hysteresis picture in which forward and reverse transformations are completely separated. At the same time, the energetic ordering of substitutions shows that epistasis imposes a reproducible bias on how one enzymatic background can be converted into the other.

\begin{figure}[H]
    \centering
    \includegraphics[width=0.82\textwidth]{Fig7.png}
    \caption[Sequence-space context of EVB trajectories]{Geometric projection of the mutational paths connecting mouse and human GPX6. Blue and orange points represent individual mutational intermediates along the stepwise substitution trajectories from mouse to human and from human to mouse, respectively. Each point is labeled by the amino acid selected at each level following the greedy algorithm criteria, showing how sequence changes accumulate across the two directional paths relative to the node~25 ancestor (green).}
    \label{fig:gpx6_triangle}
\end{figure}

\subsection{ProteinMPNN as a Secondary Test of Ancestral Sequence Recovery}

The EVB paths were selected from catalytic barriers, not from sequence probabilities. We therefore used ProteinMPNN as a separate test of which sequences are compatible with the sampled structures. Designs generated from the two sets of structures overlap in parts of sequence space but remain distinct from the reconstructed node~25 ancestor (Figures~\ref{SI-fig:proteinmpnn_design} and \ref{SI-fig:gpx6_mpnn_heatmap}). Low catalytic barriers and similarity to the ancestral sequence are therefore related but distinct properties, and the sampled structures alone are not sufficient to recover the ancestral GPX6 sequence.
